%% ============================================================
%% cap01.tex — Capítulo 1
%% Fundamentos de POO y contexto institucional
%% ============================================================

\chapter{Fundamentos de \poo\ y contexto institucional}
\label{cap:fundamentos}

\noindent
El presente capítulo sitúa la \poo\ en el conjunto más amplio de los
paradigmas de programación y establece el marco institucional dentro
del cual se desarrollará el caso de estudio a lo largo de las doce
iteraciones del libro: el sistema de información para el Estatuto General
(\acuerdo) y el Estatuto Académico de la \ud.
Puesto que el dominio del problema determina en buena medida las
decisiones de diseño orientado a objetos, se considera indispensable
que el estudiante comprenda el entramado operativo de la institución
antes de escribir la primera línea de código.
La experiencia docente confirmaa que el estudiante asimila con mayor
profundidad los conceptos de encapsulamiento y herencia cuando los aplica
sobre un dominio que le resulta familiar, habida cuenta de que la universidad
no es un sistema bancario de manual: es su propio entorno cotidiano,
con reglas, actores y tensiones que él mismo ha experimentado.
%% ─────────────────────────────────────────────
\section{Paradigmas de programación y el enfoque orientado a objetos}
\label{sec:paradigmas}
%% ─────────────────────────────────────────────

Cuatro paradigmas principales han dominado la historia de la programación
desde los años sesenta: el imperativo, el funcional, el lógico y el orientado
a objetos. La Tabla~\ref{tab:paradigmas} los contrasta en las dimensiones
que resultan más pertinentes para la discusión que se desarrollará a lo
largo del libro.

\begingroup
\small
\begin{longtable}{L{2.2cm} L{2.5cm} L{2.8cm} L{3cm} L{2.8cm}}
	% --- Configuración de encabezados y pies ---
	\caption{Comparación de los cuatro paradigmas de programación principales. Fuente: elaboración del autor con base en \textcite{Sommerville2015} y \textcite{Pressman2014}.} \label{tab:paradigmas} \\
	\toprule
	\textbf{Paradigma} & \textbf{Abstracción principal} & \textbf{Lenguajes representativos} & \textbf{Fortalezas} & \textbf{Limitaciones} \\
	\midrule
	\endfirsthead
	
	\multicolumn{5}{c}{{\bfseries \tablename\ \thetable{} -- Continuación}} \\
	\toprule
	\textbf{Paradigma} & \textbf{Abstracción principal} & \textbf{Lenguajes representativos} & \textbf{Fortalezas} & \textbf{Limitaciones} \\
	\midrule
	\endhead
	
	\midrule
	\multicolumn{5}{r}{{Continúa en la siguiente página...}} \\
	\bottomrule
	\endfoot
	
	\bottomrule
	\endlastfoot
	
	% --- Datos de la tabla ---
	Imperativo / Procedimental
	& Instrucción secuencial y procedimiento
	& C, Pascal, COBOL, Fortran
	& Cercano al hardware; eficiencia de ejecución; amplio soporte de herramientas
	& Escalabilidad limitada; acoplamiento entre datos y lógica; dificultad de mantenimiento en sistemas grandes \\
	\midrule
	
	Funcional
	& Función matemática pura; composición; inmutabilidad
	& Haskell, Erlang, Clojure, Scala, F\#
	& Razonamiento formal facilitado; ausencia de efectos secundarios; concurrencia natural
	& Curva de aprendizaje pronunciada; rendimiento en operaciones con estado mutable; ecosistema más reducido en dominios empresariales \\
	\midrule
	
	Lógico
	& Predicado, regla y unificación
	& Prolog, Datalog, Answer Set Programming
	& Expresividad en dominios de conocimiento; búsqueda automática de soluciones; adecuado para sistemas expertos
	& Rendimiento impredecible; dificultad para modelar efectos secundarios; ecosistema académico más que industrial \\
	\midrule
	
	Orientado a objetos
	& Clase, objeto, mensaje
	& Java, C++, Python, C\#, Kotlin, \textit{Rust} (parcial)
	& Modelado natural de dominios complejos; reutilización por herencia y composición; encapsulamiento del estado; patrones de diseño maduros
	& Herencia mal empleada genera acoplamiento fuerte; sobrecarga conceptual inicial; dificultad en sistemas altamente concurrentes sin cuidado explícito \\
\end{longtable}
\endgroup

\subsection*{Los cuatro pilares del paradigma orientado a objetos}

El paradigma orientado a objetos (-POO-) descansa sobre cuatro principios cuya
comprensión es condición necesaria ---aunque no suficiente--- para diseñar
sistemas mantenibles.\footnote{%
  \textcite{Gamma1994} advierten que la herencia es con frecuencia
  sobreempleada por programadores novatos. La recomendación del libro
  es preferir la composición sobre la herencia (\textit{favor composition over
  inheritance}), salvo cuando la relación «es un» sea verdaderamente estable
  en el dominio. En el dominio de la \ud, la relación entre \texttt{Docente}
  y \texttt{Persona} parece estable; la relación entre
  \texttt{DirectorDeEscuela} y \texttt{Docente} podría modelarse mejor
  como un rol.}

\textbf{Encapsulamiento.} Los datos internos de un objeto quedan ocultos
al exterior. Solo se expone lo necesario mediante una interfaz pública.
En el dominio de la \ud, la clase \texttt{Facultad} encapsula su lista
de escuelas adscritas; ningún otro objeto puede modificarla directamente,
sino solo a través de métodos definidos en la propia clase.
Así, los invariantes del dominio ---como el requisito de que una Escuela
no puede pertenecer simultáneamente a dos Facultades--- se preservan
mediante el encapsulamiento.

\textbf{Herencia.} Una clase puede derivarse de otra, heredando sus atributos
y comportamientos mientras añade o sobreescribe los que le son propios.
La clase abstracta \texttt{UnidadAcademica} reúne los atributos comunes
a Facultad, Escuela, Centro e Instituto ---nombre, código, fecha de creación,
director---; cada subclase agrega los atributos y comportamientos específicos
que la distinguen.

\textbf{Polimorfismo.} Un mismo mensaje puede producir comportamientos distintos
según el tipo del objeto receptor. El método \texttt{getTipo()} devuelve
\texttt{"Facultad"} cuando se invoca sobre un objeto de la clase
\texttt{Facultad} y \texttt{"Escuela"} cuando se invoca sobre un objeto
de la clase \texttt{Escuela}: una sola llamada, comportamientos disímiles.
La utilidad resulta evidente cuando el sistema debe generar reportes de
composición institucional sin conocer de antemano el tipo concreto de cada
unidad académica.

\textbf{Abstracción.} Consiste en identificar las características esenciales
de un objeto del dominio y descartar los detalles irrelevantes para el
propósito del sistema. El modelado de \texttt{Docente} en el sistema de
información no requiere su fecha de nacimiento ni su estado civil: requiere
su tipo de vinculación, su Escuela de adscripción y los espacios académicos
que orienta. La abstracción es, en este sentido, una decisión de diseño
contextualizada, no una operación algorítmica mecánica.

\subsection*{Pseudocódigo comparativo: paradigma imperativo vs.\ orientado a objetos}

El siguiente par de fragmentos ilustra la diferencia entre abordar el
problema de ``registrar un docente en el sistema'' desde el paradigma
imperativo y desde el paradigma orientado a objetos.

\begin{lstlisting}[style=javaStyle,
  caption={Paradigma imperativo: registrar docente mediante procedimiento y
    arreglo global.},
  label={lst:imperativo-docente}]
// Paradigma imperativo (pseudocódigo en Java-like)
String[] nombres = new String[100];
String[] vinculaciones = new String[100];
String[] escuelas = new String[100];
int totalDocentes = 0;

void registrarDocente(String nombre,
                      String vinculacion,
                      String escuela) {
    nombres[totalDocentes] = nombre;
    vinculaciones[totalDocentes] = vinculacion;
    escuelas[totalDocentes] = escuela;
    totalDocentes++;
}
\end{lstlisting}

\begin{lstlisting}[style=javaStyle,
  caption={Paradigma orientado a objetos: registrar docente encapsulando
    estado y comportamiento en la clase \texttt{Escuela}.},
  label={lst:oo-docente}]
// Paradigma orientado a objetos
public class Docente {
    private String nombre;
    private TipoDocente vinculacion;
    private Escuela escuelaAdscrita;

    public Docente(String nombre,
                   TipoDocente vinculacion,
                   Escuela escuelaAdscrita) {
        this.nombre = nombre;
        this.vinculacion = vinculacion;
        this.escuelaAdscrita = escuelaAdscrita;
    }

    public TipoDocente getVinculacion() {
        return vinculacion;
    }
}

public class Escuela extends UnidadAcademica {
    private List<Docente> docentesAdscritos = new ArrayList<>();

    public void adscribirDocente(Docente d) {
        docentesAdscritos.add(d);
    }
}
\end{lstlisting}

La diferencia no es solo sintáctica. El fragmento imperativo distribuye los
datos del docente en tres arreglos paralelos, sin ninguna garantía de
consistencia: nada impide que \texttt{totalDocentes} supere el tamaño del
arreglo, ni que una modificación en uno de los arreglos no se refleje en
los demás. El fragmento orientado a objetos agrupa datos y comportamiento;
la invariante ``un docente pertenece exactamente a una Escuela'' queda
expresada directamente en el constructor.
No es trivial.

%% ─────────────────────────────────────────────
\section{La Universidad Distrital: estructura académico-administrativa
  según el \acuerdo}
\label{sec:estructura-ud}
%% ─────────────────────────────────────────────

El \acuerdo\ del Consejo Superior Universitario redefinió la organización
académico-administrativa de la institución con un nivel de detalle normativo
que, desde la perspectiva del modelado orientado a objetos, resulta
extraordinariamente valioso: cada artículo del estatuto aporta restricciones
de dominio que se traducen directamente en invariantes de clase, precondiciones
de método y relaciones de asociación o composición.

\subsection{Facultades, escuelas, centros, institutos y CABAS}
\label{subsec:unidades-academicas}

La Tabla~\ref{tab:unidades-academicas} describe las unidades académico-administrativas
definidas en el \acuerdo, con referencia explícita a los artículos pertinentes.

\begingroup
\small
\begin{longtable}{L{2cm} L{4.5cm} C{1.6cm} L{4cm}}
	% --- Título y encabezado de la primera página ---
	\caption{Unidades académico-administrativas de la \ud\ según el \acuerdo.} \label{tab:unidades-academicas} \\
	\toprule
	\textbf{Entidad} & \textbf{Definición según el estatuto} & \textbf{Nivel\newline jerárquico} & \textbf{Artículo del \acuerdo} \\
	\midrule
	\endfirsthead
	
	% --- Encabezado para las páginas siguientes ---
	\multicolumn{4}{c}{{\bfseries \tablename\ \thetable{} -- Continuación}} \\
	\toprule
	\textbf{Entidad} & \textbf{Definición según el estatuto} & \textbf{Nivel\newline jerárquico} & \textbf{Artículo del \acuerdo} \\
	\midrule
	\endhead
	
	% --- Pie de tabla para páginas intermedias ---
	\midrule
	\multicolumn{4}{r}{{Continúa en la siguiente página...}} \\
	\bottomrule
	\endfoot
	
	% --- Pie de tabla final ---
	\bottomrule
	\endlastfoot
	
	% --- Datos de la tabla ---
	Facultad
	& Unidad académico-administrativa de primer nivel; agrupa Escuelas, Centros, Institutos y CABAS; cuenta con Consejo de Facultad y Decanatura
	& 1°
	& Arts.\ 14, 18 \\
	\midrule
	
	Escuela
	& Unidad de adscripción de los docentes; responsable de la organización, gestión y evaluación curricular; cuenta con Consejo de Escuela y Dirección de Escuela
	& 2°
	& Arts.\ 15, 18 \\
	\midrule
	
	Centro
	& Unidad con consejo y dirección propios; orientada a funciones específicas de investigación, extensión o apoyo
	& 2°
	& Art.\ 16 \\
	\midrule
	
	Instituto
	& Unidad con consejo y dirección propios; puede articular docencia, investigación y extensión en un área disciplinar particular
	& 2°
	& Art.\ 17 \\
	\midrule
	
	CABA\newline (Comité Académico Básico de Área)
	& Comité que articula docencia e investigación dentro de un área de formación; no constituye una unidad autónoma sino un mecanismo de coordinación curricular
	& 3° (coordinación)
	& Art.\ 19 \\
\end{longtable}
\endgroup

Tres aspectos merecen atención desde la perspectiva del modelado.
Primero: la Escuela no ofrece programas académicos de manera directa
---esa responsabilidad recae en las Áreas de Formación y en los propios
programas---, pero sí adscribe a todos los docentes que participan en
dichos programas. Esta asimetría genera, en el diagrama de clases, una
relación de asociación entre \texttt{Escuela} y \texttt{Docente} que
no coincide con la relación de composición entre \texttt{Facultad} y
\texttt{ProgramaAcademico}. Segundo: los Centros e Institutos tienen
consejo y dirección propios, lo cual sugiere que comparten estructura
de gobierno con las Facultades, aunque ocupen un nivel jerárquico diferente.
El patrón de diseño \textit{Composite} podría modelar esta recursividad,
según se discutirá en el Capítulo~\ref{cap:uml}. Tercero ---y esto queda
abierto como pregunta de diseño---: ¿debe el CABA modelarse como una clase
o como una relación de asociación entre \texttt{Escuela} y
\texttt{AreaDeFormacion}?

\subsection{Comunidad universitaria: docentes, estudiantes,
  personal administrativo, egresados}
\label{subsec:comunidad-universitaria}

El Artículo 8 del \acuerdo\ define la comunidad universitaria e identifica
sus componentes. La Tabla~\ref{tab:comunidad} los detalla con sus tipos
de vinculación.

\begingroup
\small
\begin{longtable}{L{2.8cm} L{4cm} L{5.5cm}}
	% --- Título y encabezado de la primera página ---
	\caption{Miembros de la comunidad universitaria según el Art.\ 8 del \acuerdo.} \label{tab:comunidad} \\
	\toprule
	\textbf{Miembro} & \textbf{Tipo de vinculación} & \textbf{Características relevantes para el modelado} \\
	\midrule
	\endfirsthead
	
	% --- Encabezado para las páginas siguientes ---
	\multicolumn{3}{c}{{\bfseries \tablename\ \thetable{} -- Continuación}} \\
	\toprule
	\textbf{Miembro} & \textbf{Tipo de vinculación} & \textbf{Características relevantes para el modelado} \\
	\midrule
	\endhead
	
	% --- Pie de tabla para páginas intermedias ---
	\midrule
	\multicolumn{3}{r}{{Continúa en la siguiente página...}} \\
	\bottomrule
	\endfoot
	
	% --- Pie de tabla final ---
	\bottomrule
	\endlastfoot
	
	% --- Datos de la tabla ---
	Docente de planta
	& Empleado público de carrera docente
	& Vinculación estable; con derechos de escalafón; adscrito a una Escuela \\
	\midrule
	
	Docente ocasional
	& Contrato de prestación de servicios especiales
	& Vinculación temporal; con derechos reducidos; adscrito a una Escuela \\
	\midrule
	
	Docente hora cátedra
	& Contrato de prestación de servicios por hora
	& Vinculación por período académico; sin continuidad garantizada \\
	\midrule
	
	Docente visitante
	& Invitado de otra institución; vinculación temporal
	& Puede no estar adscrito a una Escuela de manera formal \\
	\midrule
	
	Docente experto
	& Profesional con experiencia certificada; sin grado académico requerido
	& Vinculación especial en áreas técnico-profesionales \\
	\midrule
	
	Estudiante de pregrado
	& Matrícula vigente en programa de pregrado
	& Identificado por código de matrícula; asociado a un Programa \\
	\midrule
	
	Estudiante de posgrado
	& Matrícula vigente en programa de posgrado
	& Puede simultanear con vinculación como auxiliar de investigación \\
	\midrule
	
	Personal administrativo
	& Empleado público o trabajador oficial
	& Rol de soporte; interactúa con el sistema como usuario administrativo \\
	\midrule
	
	Egresado
	& Graduado de un programa de la \ud
	& Sin matrícula vigente; puede conservar derechos de uso de servicios \\
\end{longtable}
\endgroup

La multiplicidad de tipos de docente plantea un dilema de modelado que se
anticipó brevemente en el prefacio y que se resolverá en el
Capítulo~\ref{cap:uml}: ¿debe \texttt{TipoDocente} ser un enumerado, una
jerarquía de subclases o un conjunto de atributos booleanos? La respuesta
depende de si los tipos difieren en comportamiento ---métodos--- o solo en
datos ---atributos---. Si un docente visitante y un docente de planta
responden de manera idéntica a todos los mensajes del sistema, el enumerado
bastará; si sus comportamientos son disímiles ---por ejemplo, si la liquidación
de honorarios sigue reglas distintas según el tipo de vinculación---,
la jerarquía de subclases resulta más expresiva.

\subsection{Órganos de gobierno y toma de decisiones}
\label{subsec:organos-gobierno}

Los órganos de gobierno de la \ud\ establecidos en el Artículo 18 del
\acuerdo\ se presentan en la Tabla~\ref{tab:organos-gobierno}.

\begingroup
\small
\begin{longtable}{L{3cm} L{4.5cm} L{5.5cm}}
	% --- Título y encabezado de la primera página ---
	\caption{Órganos de gobierno de la \ud\ según el Art.\ 18 del \acuerdo.} \label{tab:organos-gobierno} \\
	\toprule
	\textbf{Órgano} & \textbf{Composición} & \textbf{Función principal} \\
	\midrule
	\endfirsthead
	
	% --- Encabezado para las páginas siguientes ---
	\multicolumn{3}{c}{{\bfseries \tablename\ \thetable{} -- Continuación}} \\
	\toprule
	\textbf{Órgano} & \textbf{Composición} & \textbf{Función principal} \\
	\midrule
	\endhead
	
	% --- Pie de tabla para páginas intermedias ---
	\midrule
	\multicolumn{3}{r}{{Continúa en la siguiente página...}} \\
	\bottomrule
	\endfoot
	
	% --- Pie de tabla final ---
	\bottomrule
	\endlastfoot
	
	% --- Datos de la tabla ---
	Consejo Superior Universitario
	& Representantes de gobierno nacional, Presidente de la República, MEN, Gobernación, comunidad académica, estudiantes, egresados, sector productivo
	& Máximo órgano de dirección y gobierno; expide el Estatuto General y los acuerdos de mayor jerarquía \\
	\midrule
	
	Consejo Académico
	& Rector, Vicerrectores, Decanos, representantes docentes y estudiantiles
	& Dirección académica; expide normas de orden académico; aprueba programas y reformas curriculares \\
	\midrule
	
	Rectoría
	& Rector, elegido por el Consejo Superior
	& Representación legal; dirección ejecutiva de la institución \\
	\midrule
	
	Consejo de Facultad
	& Decano, directores de Escuela, representantes docentes y estudiantiles
	& Dirección de la Facultad; decisiones sobre programas y docentes en el nivel de Facultad \\
	\midrule
	
	Consejo de Escuela
	& Director de Escuela, docentes adscritos, representante estudiantil
	& Gestión académica de la Escuela; asignación de carga docente \\
	\midrule
	
	Consejo de Centro / Instituto
	& Director, docentes vinculados, representante estudiantil (cuando aplica)
	& Dirección del Centro o Instituto correspondiente \\
	\midrule
	
	Consejo de Bienestar y Buen Vivir
	& Vicerrector delegado, representantes de la comunidad universitaria
	& Política de bienestar institucional \\
\end{longtable}
\endgroup

\noindent
El siguiente diagrama PlantUML sintetiza la estructura organizacional
de la \ud\ tal como se modelará en el sistema de información:

\begin{lstlisting}[style=plantumlStyle,
  caption={Diagrama de estructura organizacional de la \ud\ (PlantUML).},
  label={lst:plantuml-org}]
@startuml estructura_ud
skinparam classAttributeIconSize 0
skinparam classFontStyle bold
skinparam classBackgroundColor #F8F8FF
skinparam classBorderColor #00467F
skinparam ArrowColor #00467F

package "Universidad Distrital FJC" {
  class Universidad {
    - nombre : String
    - nit : String
    + getFacultades() : List<Facultad>
  }
  class Facultad {
    - nombre : String
    - decano : String
    + getEscuelas() : List<Escuela>
  }
  class Escuela {
    - nombre : String
    - directorEscuela : String
    + getDocentesAdscritos() : List<Docente>
  }
  class Centro {
    - nombre : String
    - director : String
  }
  class Instituto {
    - nombre : String
    - director : String
  }
  class CABA {
    - areaDeFormacion : String
    - coordinador : String
  }
  class AreaDeFormacion {
    - nombre : String
  }
  class ProgramaAcademico {
    - nombre : String
    - codigoSNIES : String
    - nivel : NivelAcademico
  }
}

Universidad "1" *-- "1..*" Facultad : contiene
Facultad "1" *-- "1..*" Escuela : agrupa
Facultad "1" o-- "0..*" Centro : contiene
Facultad "1" o-- "0..*" Instituto : contiene
Escuela "1" o-- "0..*" CABA : coordina
Facultad "1" o-- "1..*" AreaDeFormacion : organiza
AreaDeFormacion "1" --> "1..*" ProgramaAcademico : ofrece

@enduml
\end{lstlisting}

%% ─────────────────────────────────────────────
\section{El nuevo Estatuto Académico: sistema de formación,
  currículo y programas}
\label{sec:estatuto-academico}
%% ─────────────────────────────────────────────

La propuesta de Estatuto Académico de noviembre de 2025 articula el sistema
de formación universitaria en torno a cuatro componentes que, desde la
perspectiva del modelado orientado a objetos, constituyen las clases
centrales del subsistema académico. La Tabla~\ref{tab:sistema-formacion}
los describe con referencia a los artículos pertinentes.

\begingroup
\small
\begin{longtable}{L{3cm} L{8.5cm} C{1.8cm}}
	% --- Título y encabezado de la primera página ---
	\caption{Componentes del sistema de formación según el Estatuto Académico (propuesta, noviembre de 2025), Arts.\ 13--18.} \label{tab:sistema-formacion} \\
	\toprule
	\textbf{Componente} & \textbf{Descripción} & \textbf{Artículo} \\
	\midrule
	\endfirsthead
	
	% --- Encabezado para las páginas siguientes ---
	\multicolumn{3}{c}{{\bfseries \tablename\ \thetable{} -- Continuación}} \\
	\toprule
	\textbf{Componente} & \textbf{Descripción} & \textbf{Artículo} \\
	\midrule
	\endhead
	
	% --- Pie de tabla para páginas intermedias ---
	\midrule
	\multicolumn{3}{r}{{Continúa en la siguiente página...}} \\
	\bottomrule
	\endfoot
	
	% --- Pie de tabla final ---
	\bottomrule
	\endlastfoot
	
	% --- Datos de la tabla ---
	Formación universitaria
	& Proceso intencional, sistemático y contextualizado que propende al desarrollo integral del sujeto en sus dimensiones cognitiva, ética, estética, cultural y ciudadana; articulado con el Proyecto Universitario Institucional (\textsc{PUI})
	& Art.\ 13 \\
	\midrule
	
	Currículo
	& Conjunto organizado de propósitos, contenidos, estrategias pedagógicas y mecanismos de evaluación que orienta el proceso de formación en un programa académico; expresado en el Proyecto Educativo de Programa (PEP)
	& Art.\ 14 \\
	\midrule
	
	Docencia
	& Práctica pedagógica profesional mediante la cual los docentes orientan los procesos de enseñanza-aprendizaje en los espacios académicos; articulada con la investigación y la extensión
	& Art.\ 15 \\
	\midrule
	
	Programa académico
	& Unidad organizativa que ofrece un conjunto de espacios académicos orientados a la formación en un campo del conocimiento; regulado por el Decreto~1330 de 2019 y el Acuerdo~02 de 2020 del CESU; requiere registro calificado
	& Art.\ 18 \\
\end{longtable}
\endgroup

La cadena normativa PUI $\rightarrow$ PEF (Proyecto Educativo de Facultad)
$\rightarrow$ PEP (Proyecto Educativo de Programa) constituye una jerarquía
de documentos de política que el sistema de información deberá gestionar.
Cada nivel hereda los propósitos del nivel superior y los particulariza
para su unidad académica. El patrón de diseño \textit{Template Method}
podría resultar adecuado para modelar esta particularización, según se
analizará en el Capítulo~\ref{cap:java}.\footnote{%
  La cadena PUI $\rightarrow$ PEF $\rightarrow$ PEP no es un requisito
  formal de los estatutos de otras universidades colombianas. La arquitectura
  de tres niveles es específica de la \ud\ y, por lo tanto, constituye
  un argumento adicional para la elección de la institución como caso de
  estudio: el dominio es auténtico e imposible de generar mediante extrapolación
  a partir de sistemas genéricos.}

%% ─────────────────────────────────────────────
\section{Propósitos de Formación y de Aprendizaje (-PFA-) en la \ud}
\label{sec:pfa}
%% ─────────────────────────────────────────────

Los Propósitos de Formación y de Aprendizaje (-PFA-) son la herramienta
conceptual mediante la cual la \ud\ articula, en el nivel del programa y del
espacio académico, las intenciones formativas del proyecto curricular con las
estrategias de evaluación en el aula. El documento institucional de 2023
---al momento de redactar este capítulo, la versión más reciente disponible---
los organiza en tres ámbitos que la Tabla~\ref{tab:pfa-evaluacion} detalla
junto con su relación con los mecanismos de evaluación formativa.

\begingroup
\small
\begin{longtable}{L{2.5cm} L{4.5cm} L{6cm}}
	% --- Título y encabezado de la primera página ---
	\caption{Ámbitos del -PFA- y su relación con la evaluación formativa. Fuente: Documento -PFA- \ud\ (2023) y \textcite{Decreto1330_2019}.} \label{tab:pfa-evaluacion} \\
	\toprule
	\textbf{Ámbito del -PFA-} & \textbf{Descripción} & \textbf{Relación con evaluación formativa} \\
	\midrule
	\endfirsthead
	
	% --- Encabezado para las páginas siguientes ---
	\multicolumn{3}{c}{{\bfseries \tablename\ \thetable{} -- Continuación}} \\
	\toprule
	\textbf{Ámbito del -PFA-} & \textbf{Descripción} & \textbf{Relación con evaluación formativa} \\
	\midrule
	\endhead
	
	% --- Pie de tabla para páginas intermedias ---
	\midrule
	\multicolumn{3}{r}{{Continúa en la siguiente página...}} \\
	\bottomrule
	\endfoot
	
	% --- Pie de tabla final ---
	\bottomrule
	\endlastfoot
	
	% --- Datos de la tabla ---
	Ontológico
	& Se interroga por el tipo de sujeto que la formación busca contribuir a constituir; por el ser del estudiante como persona, ciudadano y profesional en devenir
	& Evaluación mediante portafolios de reflexión personal y evidencias de desarrollo ético-ciudadano; difícil de reducir a indicadores numéricos \\
	\midrule
	
	Epistemológico
	& Se ocupa del conocimiento que el estudiante debe construir, las formas de acceder a él y los criterios de validación propios del campo disciplinar
	& Evaluación mediante rúbricas de desempeño, tareas auténticas y proyectos de investigación formativa; articulado con los Resultados de Aprendizaje del Decreto~1330 de 2019 \\
	\midrule
	
	Contextual
	& Sitúa la formación en el contexto social, cultural, económico y político; propende a la pertinencia del programa frente a las necesidades del entorno
	& Evaluación mediante proyectos de impacto social, prácticas profesionales y sistematización de experiencias comunitarias \\
\end{longtable}
\endgroup

La articulación entre los -PFA- y el sistema de evaluación formal ---rúbricas,
evidencias, calificaciones--- es uno de los retos más complejos del modelado.
No es descartable que algunos ámbitos del -PFA- resistan la formalización
en clases de datos con atributos discretos: el ámbito ontológico, en
particular, parece resistirse a la representación tabular. Sin embargo,
los ámbitos epistemológico y contextual sí generan requisitos funcionales
concretos para el sistema de información, toda vez que se articulan con
los Resultados de Aprendizaje (RA) formulados mediante verbos de la taxonomía
de Bloom \autocite{MEN2020_RA}.

%% ─────────────────────────────────────────────
\section{Caso de estudio: sistema de información para el Estatuto
  General y Académico}
\label{sec:caso-estudio}
%% ─────────────────────────────────────────────

El sistema de información que actuará como hilo conductor del libro se
estructura en cinco subsistemas funcionales, cuyo alcance queda delimitado
en la lista siguiente:

\begin{enumerate}[label=\textbf{SS-\arabic*.}]
  \item \textbf{Gestión organizacional:} creación, consulta y modificación
    de Facultades, Escuelas, Centros, Institutos y CABAS; relaciones de
    pertenencia y jerarquía; gestión de directivos.
  \item \textbf{Gestión de comunidad:} registro y consulta de docentes (con
    sus cinco tipos de vinculación), estudiantes (pregrado y posgrado),
    personal administrativo y egresados; gestión de adscripciones.
  \item \textbf{Gestión académica:} programas académicos, planes de estudio,
    espacios académicos, PEP; matrícula de estudiantes; asignación de
    docentes a espacios académicos.
  \item \textbf{Gestión de evaluación y -PFA-:} Resultados de Aprendizaje,
    Propósitos de Formación y de Aprendizaje, rúbricas de evaluación,
    evidencias de aprendizaje, seguimiento del avance del estudiante.
  \item \textbf{Gestión de órganos colegiados:} registro de sesiones y
    decisiones del Consejo Superior, Consejo Académico, Consejos de Facultad,
    Escuela, Centro e Instituto; convocatorias y actas.
\end{enumerate}

\noindent
Los módulos principales del sistema, en correspondencia con los subsistemas
anteriores, son:

\begin{itemize}
  \item Módulo Organizacional (\texttt{org}): clases \texttt{Universidad},
    \texttt{Facultad}, \texttt{Escuela}, \texttt{Centro}, \texttt{Instituto},
    \texttt{CABA}, \texttt{AreaDeFormacion}.
  \item Módulo Comunidad (\texttt{comunidad}): clases \texttt{Persona},
    \texttt{Docente} (con \texttt{TipoDocente}), \texttt{Estudiante},
    \texttt{PersonalAdministrativo}, \texttt{Egresado}.
  \item Módulo Académico (\texttt{academico}): clases \texttt{ProgramaAcademico},
    \texttt{PlanDeEstudios}, \texttt{EspacioAcademico},
    \texttt{PEP}, \texttt{Matricula}.
  \item Módulo Evaluación (\texttt{evaluacion}): clases
    \texttt{ResultadoDeAprendizaje}, \texttt{PropositoDeFormacion},
    \texttt{Rubrica}, \texttt{Criterio}, \texttt{Evidencia},
    \texttt{Calificacion}.
  \item Módulo Gobierno (\texttt{gobierno}): clases
    \texttt{OrganoColegado}, \texttt{Sesion}, \texttt{Acta},
    \texttt{Acuerdo}, \texttt{Resolucion}.
\end{itemize}

%% ─────────────────────────────────────────────
\section{Metodología ágil con doce \textit{sprints}: visión general}
\label{sec:sprints-vision}
%% ─────────────────────────────────────────────

El libro adopta una metodología pedagógica de doce \textit{sprints} de
duración aproximada de una semana cada uno. La Tabla~\ref{tab:sprints-resumen}
sintetiza el foco, el capítulo del libro y los artefactos principales de
cada \textit{sprint}.

\begin{longtable}{C{1.2cm} L{3.5cm} C{1.8cm} L{5cm}}
  \toprule
  \textbf{Sprint} & \textbf{Foco temático} & \textbf{Capítulo} &
  \textbf{Artefactos principales} \\
  \midrule
  \endhead
  \bottomrule
  \caption{Resumen de los doce \textit{sprints} del libro.}
  \label{tab:sprints-resumen}
  \endlastfoot
  1 & Comprensión del dominio institucional & 1 &
    Glosario de entidades; diagrama de contexto inicial \\
  \midrule
  2 & Levantamiento de requisitos & 2 &
    Documento RF/RNF; \textit{product backlog} inicial \\
  \midrule
  3 & Validación de requisitos & 2 &
    Acta de validación; \textit{backlog} refinado con criterios de aceptación \\
  \midrule
  4 & Modelo de dominio inicial & 3 &
    Diagrama de clases: subsistemas organizacional y comunidad \\
  \midrule
  5 & Casos de uso y secuencias & 3 &
    Diagrama CU completo; diagramas de secuencia principales \\
  \midrule
  6 & Implementación \java: organización y comunidad & 4 &
    Clases \texttt{UnidadAcademica}, \texttt{Persona} y subclases en \java \\
  \midrule
  7 & Implementación \java: académico y evaluación & 4 &
    Clases \texttt{ProgramaAcademico}, \texttt{Rubrica}, \texttt{Evidencia} \\
  \midrule
  8 & Implementación \rust: \textit{structs} y \textit{traits} & 5 &
    Módulos \texttt{org} y \texttt{comunidad} en \rust; comparación con \java \\
  \midrule
  9 & Implementación \rust: académico y evaluación & 5 &
    Módulos \texttt{academico} y \texttt{evaluacion} en \rust \\
  \midrule
  10 & Pruebas unitarias e integración & 6 &
    Suite de pruebas JUnit y Cargo test; informe de cobertura \\
  \midrule
  11 & Refactorización y patrones de diseño & 6 &
    Aplicación de patrones \textit{Composite}, \textit{State},
    \textit{Observer}; métricas de calidad \\
  \midrule
  12 & Reflexión crítica: software generativo y ética & 7 &
    Informe de análisis crítico; rúbrica de evaluación de asistentes \\
\end{longtable}

La progresión sigue un arco deliberado: del análisis del dominio (Sprints~1--3)
al modelado (Sprints~4--5), a la implementación (Sprints~6--9), a la
verificación y mejora (Sprints~10--11) y, finalmente, a la reflexión
crítica sobre el proceso (Sprint~12). No es una secuencia rígida; el
estudiante deberá revisitar decisiones de análisis cuando la implementación
revele ambigüedades en el modelo.

%% ─────────────────────────────────────────────
\section{Comparativa \java\ vs.\ \rust\ para el dominio académico}
\label{sec:java-rust-intro}
%% ─────────────────────────────────────────────

La elección de \java\ y \rust\ como lenguajes de implementación no es
arbitraria. La Tabla~\ref{tab:java-rust-comparativa} los contrasta en las
dimensiones más relevantes para el dominio académico-administrativo de la \ud.

\begingroup
\small
\begin{longtable}{L{2.8cm} L{5.5cm} L{5.5cm}}
	% --- Título y encabezado de la primera página ---
	\caption{Comparativa \java\ vs.\ \rust\ en dimensiones relevantes para el dominio académico. Fuentes: \textcite{Bloch2018} y \textcite{Klabnik2023}.} \label{tab:java-rust-comparativa} \\
	\toprule
	\textbf{Dimensión} & \textbf{\java} & \textbf{\rust} \\
	\midrule
	\endfirsthead
	
	% --- Encabezado para las páginas siguientes ---
	\multicolumn{3}{c}{{\bfseries \tablename\ \thetable{} -- Continuación}} \\
	\toprule
	\textbf{Dimensión} & \textbf{\java} & \textbf{\rust} \\
	\midrule
	\endhead
	
	% --- Pie de tabla para páginas intermedias ---
	\midrule
	\multicolumn{3}{r}{{Continúa en la siguiente página...}} \\
	\bottomrule
	\endfoot
	
	% --- Pie de tabla final ---
	\bottomrule
	\endlastfoot
	
	% --- Datos de la tabla ---
	Tipado
	& Estático; tipado nominal; genéricos con borrado de tipos en tiempo de ejecución; \textit{wildcards} (\texttt{? extends}, \texttt{? super})
	& Estático; tipado nominal y estructural; genéricos monomorfizados; sistema de tipos más expresivo con \textit{traits} \\
	\midrule
	
	Manejo de memoria
	& Recolector de basura (\textit{garbage collector}); sin gestión manual; pausa del GC puede afectar latencia
	& Sin GC; sistema de \textit{ownership} con comprobación en tiempo de compilación; memoria liberada de forma determinista \\
	\midrule
	
	Concurrencia
	& Hilos con monitores (\texttt{synchronized}); \textit{java.util.concurrent}; \texttt{CompletableFuture}; riesgo de condiciones de carrera
	& Concurrencia garantizada por el compilador (\textit{fearless concurrency}); \texttt{Arc}, \texttt{Mutex}; sin carreras de datos posibles en código seguro \\
	\midrule
	
	Ecosistema -GUI-
	& JavaFX, Swing, SWT; maduros y estables; integración con herramientas de construcción (Maven, Gradle)
	& Ecosistema -GUI- menos maduro; opciones: \textit{egui}, \textit{Tauri}; preferible para \textit{backend} o servicios web \\
	\midrule
	
	Curva de aprendizaje
	& Moderada para fundamentos de -POO-; la biblioteca estándar es amplia y la documentación abundante
	& Pronunciada; el sistema de \textit{ownership} requiere cambio de modelo mental; recomendado como segundo lenguaje tras \java \\
\end{longtable}
\endgroup

Los fragmentos de código siguientes ilustran la clase \texttt{Facultad}
en ambos lenguajes:

\begin{lstlisting}[style=javaStyle,
  caption={Clase \texttt{Facultad} en \java\ (versión simplificada).},
  label={lst:java-facultad}]
import java.util.ArrayList;
import java.util.Collections;
import java.util.List;

public class Facultad extends UnidadAcademica {
    private String decano;
    private List<Escuela> escuelas;
    private List<ProgramaAcademico> programas;

    public Facultad(String nombre, String codigo, String decano) {
        super(nombre, codigo);
        this.decano = decano;
        this.escuelas = new ArrayList<>();
        this.programas = new ArrayList<>();
    }

    public void agregarEscuela(Escuela e) {
        if (e == null) throw new IllegalArgumentException(
            "La escuela no puede ser nula.");
        escuelas.add(e);
    }

    public List<Escuela> getEscuelas() {
        return Collections.unmodifiableList(escuelas);
    }

    @Override
    public String getTipo() { return "Facultad"; }

    public String getDecano() { return decano; }
}
\end{lstlisting}

\begin{lstlisting}[style=rustStyle,
  caption={\textit{Struct} \texttt{Facultad} en \rust\ con
    implementación de métodos.},
  label={lst:rust-facultad}]
#[derive(Debug, Clone)]
pub struct Facultad {
    pub nombre: String,
    pub codigo: String,
    pub decano: String,
    escuelas: Vec<Escuela>,
    programas: Vec<ProgramaAcademico>,
}

impl Facultad {
    pub fn new(nombre: &str,
               codigo: &str,
               decano: &str) -> Self {
        Facultad {
            nombre: nombre.to_string(),
            codigo: codigo.to_string(),
            decano: decano.to_string(),
            escuelas: Vec::new(),
            programas: Vec::new(),
        }
    }

    pub fn agregar_escuela(&mut self, e: Escuela) {
        self.escuelas.push(e);
    }

    pub fn escuelas(&self) -> &[Escuela] {
        &self.escuelas
    }

    pub fn tipo(&self) -> &'static str { "Facultad" }
}
\end{lstlisting}

Persiste, sin embargo, la duda de si \rust\ es apropiado como lenguaje
introductorio para estudiantes que cursan por primera vez la asignatura
de -POO- en la \ud. La experiencia observada entre 2019 y 2025 en la
Facultad de Ingeniería sugiere que el sistema de \textit{ownership} puede
suscitar más confusión que iluminación cuando el estudiante no ha consolidado
aún los conceptos de encapsulamiento y herencia. Por esta razón, el libro
presenta \java\ como lenguaje primario y \rust\ como contraste enriquecedor,
no como sustituto.

%% ─────────────────────────────────────────────
\section{Reflexión sobre software generativo en el aprendizaje
  de la programación}
\label{sec:ia-reflexion-intro}
%% ─────────────────────────────────────────────

La pregunta que atraviesa todo el libro no admite respuesta sencilla.
Los asistentes computacionales generativos ---disponibles de forma
prácticamente universal entre los estudiantes de ingeniería al momento
de redactar este texto, en marzo de 2026--- producen código \java\ y
diagramas -UML- con rapidez y con un nivel de corrección sintáctica que,
en el peor de los casos, supera el de muchos estudiantes de primer año.
Ello plantea un escollo pedagógico genuino: si el estudiante puede obtener
en segundos la clase \texttt{Facultad} con sus métodos de acceso, ¿qué
valor agrega el proceso de diseño manual?

El marco de análisis crítico adoptado en este libro propone cuatro
categorías para evaluar el código generado por asistentes computacionales:

\begin{enumerate}[label=\textbf{C\arabic*.}]
  \item \textbf{Coherencia con el dominio:} ¿el código generado refleja
    las restricciones del estatuto, o las ignora en favor de patrones
    genéricos? Un asistente que modela \texttt{Facultad} con una lista
    de \texttt{Docente} ---en lugar de una lista de \texttt{Escuela},
    ya que los docentes se adscriben a Escuelas, no a Facultades---
    evidencia una comprensión superficial del dominio.
  \item \textbf{Corrección de las relaciones:} ¿el código distingue
    entre composición, agregación y asociación simple? La diferencia
    entre contener una lista de \texttt{Escuela} (composición) y
    referenciar un \texttt{Director} (asociación) tiene implicaciones
    directas sobre la gestión del ciclo de vida de los objetos.
  \item \textbf{Aplicabilidad de principios -POO-:} ¿el código respeta
    el encapsulamiento? ¿Emplea herencia donde corresponde y composición
    donde es más adecuada? ¿Los métodos tienen nombres coherentes con el
    vocabulario del dominio?
  \item \textbf{Mantenibilidad a largo plazo:} ¿el código podría
    incorporarse en un sistema real de la \ud\ sin refactorización
    sustancial? ¿Está alineado con los principios \textsc{SOLID}
    \autocite{Martin2017}?
\end{enumerate}

Este marco se aplicará en cada uno de los doce \textit{sprints} como
parte de la actividad de reflexión. Queda por determinar si la aplicación
sistemática del marco desarrolla, en los estudiantes, una actitud genuinamente
crítica frente al software generativo o si, con el tiempo, se convierte
en un trámite formal sin incidencia real en las decisiones de diseño.\footnote{%
  La preocupación no es abstracta. En los semestres 2024-I y 2024-II,
  varios estudiantes de Ingeniería de Sistemas de la \ud\ entregaron
  código generado con mínimas modificaciones sin haber examinado si
  las clases producidas eran coherentes con el dominio del problema.
  La detección de este comportamiento llevó al autor a incorporar la
  actividad comparativa como componente de evaluación.}

%% ─────────────────────────────────────────────
%% SPRINT 1
%% ─────────────────────────────────────────────
\section*{Sprint 1: Comprensión del dominio institucional}
\label{sprint:1}
\addcontentsline{toc}{section}{Sprint 1: Comprensión del dominio institucional}

\begin{tcolorbox}[sprintBox, title={Sprint 1 — Comprensión del dominio institucional}]
  \textbf{Duración estimada:} 1 semana\\
  \textbf{Objetivo:} Identificar y caracterizar las entidades del dominio
  del \acuerdo\ y del Estatuto Académico como paso previo al modelado orientado a objetos.\\
  \textbf{Entregables:} (1) Glosario de entidades; (2) Diagrama de contexto
  inicial; (3) Tabla comparativa propia vs.\ asistente computacional.
\end{tcolorbox}

\bigskip

\subsection*{Tabla del Sprint 1}

\begin{longtable}{L{3.5cm} L{4.5cm} L{4.5cm}}
  \sprintcabecera
  \endhead
  \bottomrule
  \caption{Entregables, características del programa y criterios de la
    rúbrica del Sprint 1.}
  \label{tab:sprint1}
  \endlastfoot
  Glosario de entidades institucionales (Facultad, Escuela, Centro,
  Instituto, CABA, Docente, Estudiante, órganos de gobierno) con
  definición, nivel jerárquico y referencia al artículo del \acuerdo
  & Reconocer las entidades del dominio del \acuerdo\ y sus relaciones;
    aplicar técnicas de análisis de documentos normativos
  & El estudiante identifica correctamente al menos 15 entidades del
    dominio, las define con vocabulario coherente con el estatuto e
    infiere 5 relaciones entre ellas con argumentación normativa \\
  \midrule
  Diagrama de contexto del sistema (versión inicial, representación
  libre, no -UML- formal): actores internos, actores externos, flujos
  de información principales
  & Delimitar el alcance del sistema de información frente a sistemas
    externos (nómina, biblioteca, MEN, plataformas de \textit{e-learning})
    y actores institucionales
  & El diagrama distingue con precisión los actores internos, los actores
    externos y los flujos de información principales, con justificación
    basada en el \acuerdo \\
  \midrule
  Tabla comparativa propia vs.\ asistente computacional generativo
  (mínimo 10 filas): Entidad / Propuesta propia / Propuesta del
  asistente / Diferencias / Valoración crítica
  & Desarrollar pensamiento crítico frente al diseño generado
    automáticamente; aplicar las cuatro categorías del marco de análisis
    (Sección~\ref{sec:ia-reflexion-intro})
  & El estudiante identifica al menos 3 diferencias conceptuales
    significativas y las argumenta con base en el dominio institucional;
    evidencia comprensión de las limitaciones del asistente frente a
    dominios normativos específicos \\
  \midrule
  Bitácora de proceso (al menos 3 entradas con fecha y hora previas
  a la consulta del asistente computacional)
  & Documentar el proceso de construcción de conocimiento; distinguir
    el aprendizaje mediado por el análisis directo del estatuto del
    mediado por el asistente
  & La bitácora evidencia proceso de elaboración propio, con argumentación
    y dudas registradas, anteriores en el tiempo a la consulta del
    asistente computacional \\
\end{longtable}

\bigskip

%% ─── Tabla de alineación completa ────────────────────────
\subsection*{Tabla de alineación: libro, syllabus y resultados de aprendizaje}

\begin{longtable}{L{2.4cm} L{3.2cm} L{4cm} L{3.5cm}}
  \toprule
  \textbf{Unidad /\newline Capítulo} & \textbf{Temas del syllabus} &
  \textbf{Resultados de aprendizaje} & \textbf{Actividades /\newline Evidencias} \\
  \midrule
  \endhead
  \bottomrule
  \caption{Tabla de alineación completa: unidades del libro, syllabus,
    resultados de aprendizaje y actividades/evidencias.}
  \label{tab:alineacion}
  \endlastfoot
  Cap.~1 --- Fundamentos de -POO- y contexto institucional
  & Paradigmas de programación; estructura académica de la \ud;
    encapsulamiento, herencia, polimorfismo, abstracción; introducción
    a \java\ y \rust
  & El estudiante caracteriza los cuatro paradigmas de programación y
    distingue el modelo orientado a objetos; identifica las entidades
    del dominio institucional (\acuerdo)
  & Glosario de entidades; diagrama de contexto; tabla comparativa
    asistente/propio (Sprint~1) \\
  \midrule
  Cap.~2 --- Requisitos del sistema de información
  & Análisis de requisitos funcionales y no funcionales; historias de
    usuario; actores del sistema; priorización MoSCoW; modelo de contexto
  & El estudiante formula al menos 20 requisitos funcionales coherentes
    con el dominio normativo; distingue RF de RNF; formula historias de
    usuario con criterios de aceptación verificables
  & Documento RF/RNF; \textit{product backlog}; acta de validación;
    actividades Sprints~2 y 3 \\
  \midrule
  Cap.~3 --- Modelado -UML-
  & Diagramas de clases, casos de uso, secuencia y estado; \plantuml\
    como herramienta; principios de diseño orientado a objetos
  & El estudiante construye diagramas -UML- coherentes con el dominio;
    justifica decisiones de herencia, composición y asociación; identifica
    patrones de diseño aplicables
  & Diagramas -UML- en \plantuml; análisis crítico de modelos generados
    por asistente; actividades Sprints~4 y 5 \\
  \midrule
  Cap.~4 --- Implementación en \java
  & Clases, interfaces, herencia, polimorfismo, colecciones, excepciones;
    principios SOLID; pruebas con JUnit
  & El estudiante implementa el modelo de dominio en \java\ respetando
    los principios SOLID; aplica al menos dos patrones de diseño del
    catálogo Gang of Four \autocite{Gamma1994}
  & Código fuente en \java; informe de pruebas JUnit; actividades
    Sprints~6 y 7 \\
  \midrule
  Cap.~5 --- Implementación en \rust
  & \textit{Structs}, \textit{enums}, \textit{traits}, \textit{ownership},
    \textit{lifetimes}, módulos; comparación con \java
  & El estudiante implementa el modelo de dominio en \rust\ y explica
    las diferencias de diseño respecto de la implementación en \java;
    comprende el sistema de \textit{ownership}
  & Código fuente en \rust; informe comparativo \java/\rust; actividades
    Sprints~8 y 9 \\
  \midrule
  Cap.~6 --- Pruebas y calidad
  & Pruebas unitarias, de integración y de sistema; cobertura; métricas
    de calidad; refactorización; deuda técnica
  & El estudiante diseña una suite de pruebas con cobertura superior al
    80\%; identifica y aplica al menos una refactorización para mejorar
    la mantenibilidad del código
  & Suite de pruebas; informe de cobertura; informe de refactorización;
    actividades Sprints~10 y 11 \\
  \midrule
  Cap.~7 --- Software generativo y ética en la ingeniería
  & Marco de análisis crítico de código generado; integridad académica;
    ética del software; responsabilidad profesional
  & El estudiante evalúa sistemáticamente el código generado por
    asistentes computacionales mediante el marco de cuatro categorías;
    articula una posición argumentada sobre el uso responsable de estas
    herramientas
  & Informe de análisis crítico; presentación final; actividad Sprint~12 \\
\end{longtable}

\bigskip

%% ─── Actividad IA ────────────────────────────
\subsection*{Actividad de reflexión: software generativo y diseño orientado a objetos}
\label{act:ia-sprint1}

\begin{tcolorbox}[actividadIA,
  title={Actividad de reflexión (Sprint~1) —
         Software generativo y diseño orientado a objetos}]

\textbf{Instrucción general.}
Antes de consultar cualquier asistente computacional generativo,
el estudiante debe elaborar de manera individual una propuesta inicial
de las entidades del dominio institucional de la \ud\ y de las relaciones
entre ellas, a partir únicamente de la lectura directa del \acuerdo.
La propuesta se registra en la bitácora de proceso con fecha y hora.
Seguidamente, se solicita al asistente computacional la misma tarea
---identificar las entidades del sistema de información para la \ud\
según su Estatuto General--- y se comparan ambas propuestas de manera
sistemática.

\bigskip
\textbf{Preguntas orientadoras:}
\begin{enumerate}[label=\arabic*.]
  \item ¿Qué entidades identificó el asistente computacional que usted
        no había considerado? ¿Son pertinentes para el dominio del estatuto?
        Argumente con referencia al artículo específico del \acuerdo.
  \item ¿Qué entidades identificó usted que el asistente omitió o
        denominó de manera disímil? ¿Cuál denominación resulta más
        coherente con el texto del \acuerdo?
  \item ¿El asistente reproduce errores conceptuales sobre la estructura
        de la \ud? Identifique al menos dos (por ejemplo: ¿adscribe
        docentes a Facultades en lugar de a Escuelas?) y explique su
        origen probable.
  \item ¿En qué medida la propuesta del asistente refleja el dominio
        real de la \ud, o en cambio proyecta un sistema universitario
        genérico similar a los modelos de manual?
\end{enumerate}

\bigskip
\textbf{Producto esperado:} Tabla comparativa (mínimo 10 filas) con
columnas: Entidad / Propuesta propia / Propuesta del asistente /
Diferencias / Valoración crítica (categoría C1--C4 del marco de análisis).

\bigskip
\textbf{Rúbrica de evaluación del Sprint~1 (actividad de reflexión):}

\begin{tabular}{L{3.5cm} L{2.5cm} L{2.5cm} L{2.5cm}}
  \toprule
  \textbf{Criterio} & \textbf{Insuficiente (1)} &
  \textbf{Aceptable (2--3)} & \textbf{Sobresaliente (4--5)} \\
  \midrule
  Identificación de entidades del dominio
  & Menos de 8 entidades identificadas correctamente
  & Entre 8 y 12 entidades con denominación coherente con el estatuto
  & 15 o más entidades, todas con denominación y caracterización
    respaldada en el texto del \acuerdo \\
  \midrule
  Calidad del análisis crítico
  & El estudiante no identifica diferencias significativas o las
    diferencias son superficiales
  & El estudiante identifica 2--3 diferencias con argumentación básica
  & El estudiante identifica 4 o más diferencias significativas con
    argumentación respaldada en el dominio institucional y referencia
    a los artículos del estatuto \\
  \midrule
  Proceso de elaboración (bitácora)
  & No se evidencia proceso previo a la consulta del asistente
  & Se evidencia proceso parcial; la propuesta propia antecede a la
    consulta pero no está documentada con detalle
  & Se documenta con detalle el proceso de elaboración propio, con
    fecha y hora anteriores a la consulta; la bitácora registra dudas,
    correcciones y reflexiones del proceso \\
  \midrule
  Coherencia con el marco de análisis crítico
  & No se aplica el marco de cuatro categorías
  & Se aplica el marco de manera parcial (1--2 categorías)
  & Se aplica el marco completo (C1--C4) con ejemplos concretos y
    argumentación referenciada al dominio institucional \\
  \bottomrule
\end{tabular}

\end{tcolorbox}
